This paper presents a method for tracking the three-dimensional position and
orientation of a consumer-grade inertial measurement unit (IMU) using
quaternion-based sensor fusion. Raw accelerometer and gyroscope readings are
expressed in the sensor's body frame, which rotates freely with the device. To
recover motion in a fixed world frame, we first estimate the sensor's attitude
by means of a complementary filter---a linear combination of a low-pass
filtered accelerometer signal (stable in the long term but sensitive to
vibration) and a high-pass filtered gyroscope integral (accurate in the short
term but subject to drift). The resulting orientation, encoded as a unit
quaternion, is used at every time step to rotate the measured acceleration into
the world frame and subtract gravity. The gravity-free world-frame acceleration
is then double-integrated to yield a three-dimensional position trajectory
starting from the origin. We describe the full signal-processing pipeline, from
raw sensor acquisition to position output, detail the mathematical formulation
of quaternion rotations and the complementary filter, and discuss practical
considerations such as drift mitigation and numerical stability.

\vspace{0.3em}
\noindent\textbf{Keywords:} sensor fusion, quaternions, complementary filter,
inertial navigation, accelerometer, gyroscope, attitude estimation, position
tracking.
